\chapter{Amniocentesis (Amniotic Fluid Test)}

\section*{What Is This Test?}
Amniocentesis is a prenatal diagnostic procedure in which a small sample of amniotic fluid is collected through the abdomen using ultrasound guidance. Fetal cells and fluid can be tested for selected genetic, chromosomal, biochemical, or infectious conditions.

\section*{Alternative Names}
\begin{itemize}
  \item Amniotic Fluid Analysis
  \item Prenatal Diagnostic Testing by Amniocentesis
  \item Amnio
\end{itemize}

\section*{Principle}
The fluid contains cells and molecules originating from the fetus and surrounding tissues. Testing may include chromosome analysis, chromosomal microarray, targeted molecular testing, alpha-fetoprotein, or infection studies according to the clinical question.

\section*{Why This Test Is Ordered}
It may be offered after an abnormal screening result or ultrasound finding, when there is a family or pregnancy history of a genetic condition, or when a specific fetal infection or other condition needs evaluation. Diagnostic testing options should be discussed with all pregnant patients, regardless of age or risk.

\section*{Primary vs Secondary Test}
Amniocentesis is a diagnostic test, not a general screening test. Cell-free DNA and other prenatal screens estimate risk; they do not replace diagnostic testing when a definitive answer is needed.

\section*{How This Test Is Performed}
After ultrasound identifies the fetus, placenta, and a safe fluid pocket, a thin needle is passed through the abdominal wall into the amniotic sac and a small amount of fluid is withdrawn. The procedure is commonly performed between 15 and 20 weeks, though timing depends on the indication.

\section*{What Preparation Is Needed}
Follow the obstetric team's instructions. Bring the indication, prior screening and ultrasound results, blood type information, and any genetic counseling documentation. Ask about medication and Rh status before the procedure.

\section*{Contraindications and Precautions}
The decision is voluntary and should follow informed counseling about benefits, limitations, alternatives, and the small risks of bleeding, infection, fluid leakage, and pregnancy loss. Results apply only to the conditions tested and do not exclude every birth defect.

\section*{Understanding Results}
A normal result means that the tested conditions were not detected by the methods used. An abnormal or uncertain result requires genetic counseling and specialist interpretation; the clinical significance may depend on the specific variant and laboratory method.

\section*{Follow-up Tests}
Follow-up may include genetic counseling, targeted parental testing, detailed ultrasound, fetal echocardiography, repeat or additional molecular testing, or consultation with maternal-fetal medicine and pediatric specialists.

\section*{References}
\begin{enumerate}
  \item MedlinePlus. Amniocentesis (Amniotic Fluid Test).
  \item American College of Obstetricians and Gynecologists. Amniocentesis.
  \item American College of Obstetricians and Gynecologists. Prenatal Genetic Screening Tests.
\end{enumerate}

\clearpage
\section*{Understanding Results and Follow-up}
The laboratory report should identify the specimen, method, units, reference interval or decision limit, and whether the result is positive, negative, abnormal, or inconclusive. Reference intervals vary with age, sex, pregnancy, specimen type, assay, and laboratory; a result outside a range is not automatically a diagnosis, and a result within a range does not always exclude disease. Discuss unexpected results with the ordering clinician, who can decide whether repeat or confirmatory testing, treatment, or referral is appropriate. Do not change medicines or delay urgent care based on an isolated result.


\section*{Regulatory Status and Authorization Date}
This chapter describes a test category, not a specific commercial product. FDA, CDSCO, EU IVDR, and MHRA approval or registration dates therefore cannot be assigned without the assay manufacturer, product name, instrument, and intended use. For laboratory-developed tests, an individual agency approval date may not exist. See the Preface for the interpretation of regulatory status.

\clearpage
